%% Everything is owned by Aditya Teraiya. Unauthorized use is prohibited.
\documentclass[11pt, a4paper]{article}
\usepackage[utf8]{inputenc}
\usepackage[margin=1in, top=1.3in]{geometry} % Adjusted top margin 
% --- Core Layout & Structure ---
\usepackage{fancyhdr, titlesec, multicol}

% --- Advanced Mathematics ---
\usepackage{amsmath, amssymb, amsfonts, mathtools, amsthm}

% --- Graphics, Colors & Custom Boxes ---
\usepackage{graphicx, xcolor, tcolorbox, tikz}

% --- Text Formatting & Lists ---
\usepackage{enumitem, url, hyperref}
\usepackage{bookmark}

% --- Tables, Figures & Captions ---
\usepackage{booktabs, array, float, caption, subcaption}

% --- Code Snippets & Citations ---
\usepackage{lipsum, listings}

% Custom Box Definition as requested
\newtcolorbox{box1}{
  colback=purple!10!white,
  colframe=purple!75!black,
  boxrule=0pt,
  leftrule=3pt,
  rightrule=0pt,
  toprule=0pt,
  bottomrule=0pt,
  sharp corners,
  left=4pt,right=4pt,top=4pt,bottom=4pt
}

\newtcolorbox{box2}{
  colback=blue!10!white,
  colframe=blue!75!black,
  boxrule=0pt,
  leftrule=3pt,
  rightrule=0pt,
  toprule=0pt,
  bottomrule=0pt,
  sharp corners,
  left=4pt,right=4pt,top=4pt,bottom=4pt
}

% Header Layout Configuration
\pagestyle{fancy}
\fancyhf{} % Clear default header and footer settings
\renewcommand{\headrulewidth}{0.5pt} % Adds a neat horizontal divider under your header metadata

% Top Left Margin Content
\fancyhead[L]{%
  \textbf{Aditya Teraiya} \\
  \textbf{\href{mailto: bmat2546@isibang.ac.in}{bmat2546}} \\
   Indian Statistical Institute, Bangalore.
}

% Top Right Margin Content
\fancyhead[R]{%
  \textbf{Subject:} $\mathbb{R}$eal Analysis \\
  \textbf{Class Homework(s)} \\
  \textbf{Date:} \today
}
\newlist{Week}{enumerate}{1}
\setlist[Week,1]{label=\textbf{Week \arabic*:}, leftmargin=*}

% Increase header height so the three rows of text do not collide with the page boundaries
\setlength{\headheight}{42pt}

\begin{document}

\begin{Week}
\item  \begin{enumerate}
    \item Each of the four cards shown below has a number on one side and a letter on the other side. What is the minimum number of cards that must be turned in order to prove the correctness of this statement: "If a card has a vowel on one side, then it has an even number on the other side?" 
     $$\boxed{A} \boxed{C} \boxed{4} \boxed{7}$$
    \begin{proof}
        We only need to check card \boxed{A} and \boxed{7}. We check \boxed{A} because if that is false, the claim doesn't hold. Also note the contrapositive of \textit{If one side has vowel then there's even number on the other side} will be \textit{If one side has odd number then there's a consonant on the other side}. So we check \boxed{7} as well. The statement says nothing when one side is a consonant so we ignore that. The \boxed{4} card will work either way so also ignored.
    \end{proof}
    \item Is there a explicit bijection from $\mathbb{N} \rightarrow \mathbb{Q}$ or other way around?
    \begin{proof}
        Yes. Consider f:
    \end{proof}
\end{enumerate}
\item  Let $f: \mathbb{R} \to \mathbb{R}$ be a continuous function. Prove that if $f$ is bounded on every closed interval, then $f$ is bounded on $\mathbb{R}$.
\item  Let $f: \mathbb{R} \to \mathbb{R}$ be a continuous function. Prove that if $f$ is bounded on every closed interval, then $f$ is bounded on $\mathbb{R}$.
\end{Week}
\end{document}
